% !TeX root = booklet-01-preview.tex
\section*{هدف و روش مطالعه}
این جزوه برای دانشجویی تنظیم می‌شود که درس الکترونیک دیجیتال را نگذرانده است.
پیش‌نیازهای مدار، منطق و ریاضی در محل مناسب آموزش داده می‌شوند؛ بنابراین آشنایی قبلی
با تحلیل ترانزیستور، مشتق یا انتگرال فرض نمی‌شود. هدف، رسیدن از شناخت کمیت‌های مدار
به فهم روابط جریان و ظرفیت \lr{MOSFET} است.

\textbf{وضعیت این نسخه:} فقط فهرست تفصیلی پیشنهادی و یک نمونه آموزشی نوشته شده است.
عنوان‌های فهرست، به معنای تکمیل آن درس‌ها نیستند. نمونه، میزان توضیح و شیوه ارائه را
نشان می‌دهد؛ آموزش کامل پیش‌نیازها در نسخه نهایی گسترده‌تر خواهد بود.

دامنه جزوه اول، کل فصل ۱ و بخش‌های \secnum{۲.۱} تا \secnum{۲.۳} کتاب است \cite{weste}.
اسلایدهای \lr{00} و \lr{01} راهنمای ترتیب و تاکید درس‌اند \cite{intro,mos}؛
مطالب کتاب که در اسلایدها کوتاه شده‌اند نیز با ارجاع روشن پوشش داده می‌شوند.
مثال‌های تکمیلی با عنوان «مثال آموزشی» از مثال‌های کتاب متمایز خواهند بود.

\section*{فهرست تفصیلی پیشنهادی جزوه اول}
\subsection*{بخش آغازین: پیش‌نیازهای ضروری}
\begin{enumerate}
\item \textbf{کمیت‌ها و واحدها:} بار، جریان، ولتاژ، انرژی و توان؛ جهت مرجع؛
\lr{Ground} و اختلاف پتانسیل؛ واحدهای \lr{SI} و پیشوندها.
\item \textbf{ریاضی محاسبات مدار:} کسر و نسبت، توان و نماد علمی، تبدیل واحد،
جابه‌جایی جمله‌ها، پرانتز، معادله خطی و درجه‌دو؛ انتخاب جواب سازگار با مسئله.
\item \textbf{مدارهای پایه:} مقاومت و \lr{Ohm's Law}، مدار باز و اتصال کوتاه،
\lr{Node} و \lr{Branch}؛ اتصال \lr{Series} و \lr{Parallel}؛
\lr{Kirchhoff's Current Law (KCL)} و \lr{Kirchhoff's Voltage Law (KVL)}.
\item \textbf{خازن و زمان:} بار ذخیره‌شده، رابطه بار و ولتاژ، شارژ و دشارژ؛
تفاوت حالت ماندگار و گذرا؛ مقدمه نرخ تغییر و شیب با نمودار.
\item \textbf{منطق دیجیتال:} تفاوت \lr{Analog} و \lr{Digital}؛
سطح ولتاژ و مقدار منطقی؛ \lr{Truth Table} و گیت‌های
\lr{NOT, AND, OR, NAND, NOR, XOR, XNOR}.
\item \textbf{جبر و ساختار منطق:} \lr{Boolean Algebra} و \lr{De Morgan's Laws}؛
ساخت گیت از \lr{Switch}؛ تفاوت \lr{Combinational} و \lr{Sequential}؛
مفهوم حافظه و \lr{Clock}، با مثال ساده.
\end{enumerate}

\clearpage
\subsection*{فصل اول: از ترانزیستور تا فرایند طراحی}
ترتیب زیر تمام بخش‌های \secnum{۱.۱} تا \secnum{۱.۱۲} را پوشش می‌دهد؛ مرزها با فهرست کتاب تطبیق
داده شده‌اند \cite{weste}. هر درس شامل توضیح، مثال و پرسش با پاسخ خواهد بود.
\begin{enumerate}
\item \textbf{بخش‌های \secnum{۱.۱} و \secnum{۱.۲}:} تاریخچه \lr{IC} و \lr{VLSI}،
\lr{Moore's Law} و مسیر کلی طراحی؛ خواندن رشد نمایی و نمودار نیمه‌لگاریتمی.
آمار تاریخی منبع، با تاریخ خود آن گزارش می‌شود.
\item \textbf{بخش \secnum{۱.۳}:} ساختار ابتدایی \lr{MOS}، چهار پایانه،
کنترل مسیر جریان با ولتاژ \lr{Gate}، مقایسه \lr{nMOS} و \lr{pMOS} و مدل \lr{Switch}.
\item \textbf{بخش \secnum{۱.۴}:} طراحی \lr{Inverter}، \lr{NAND} و \lr{NOR}؛
شبکه‌های \lr{Pull-up} و \lr{Pull-down} و گیت‌های \lr{Compound}؛
\lr{Pass Transistor}، \lr{Transmission Gate}، \lr{Tristate}،
\lr{Multiplexer} و مدارهای \lr{Sequential}.
\item \textbf{بخش \secnum{۱.۵}:} مقطع \lr{Inverter}، مراحل \lr{Fabrication}،
\lr{Layout Design Rules}، \lr{Gate Layout} و \lr{Stick Diagram}؛
ارتباط نمای مداری و فیزیکی با شکل‌های خوانا.
\item \textbf{بخش \secnum{۱.۶}:} \lr{Design Partitioning}، \lr{Abstraction}،
\lr{Hierarchy}، \lr{Regularity} و \lr{Modularity}؛
دامنه‌های \lr{Behavioral}، \lr{Structural} و \lr{Physical}.
\item \textbf{بخش \secnum{۱.۷}:} معرفی مقدمات لازم برای مثال \lr{MIPS}،
تفاوت \lr{Architecture} و \lr{Microarchitecture} و اجرای چندچرخه‌ای؛
ساختار پردازنده پیش از تحلیل مثال کتاب معرفی می‌شود.
\item \textbf{بخش \secnum{۱.۸}:} \lr{Logic Design}، رابط‌های سطح بالا، نمودار بلوکی،
سلسله‌مراتب و \lr{Hardware Description Language (HDL)}؛
خواندن مثال‌های کتاب بدون فرض آشنایی قبلی با کدنویسی سخت‌افزار.
\item \textbf{بخش \secnum{۱.۹}:} \lr{Circuit Design}؛ مقدمه معیارهای
\lr{Delay}، \lr{Power} و \lr{Area} در حد همین فصل و پیوند آن‌ها با ترانزیستور.
\item \textbf{بخش \secnum{۱.۱۰}:} \lr{Floorplanning}، \lr{Standard Cell}،
\lr{Pitch Matching}، \lr{Slice Plan}، \lr{Array} و برآورد مساحت.
\item \textbf{بخش‌های \secnum{۱.۱۱} و \secnum{۱.۱۲}:} هدف و سطوح \lr{Verification}؛
\lr{Fabrication}، \lr{Packaging} و \lr{Testing} و تفاوت نقش آن‌ها.
\end{enumerate}

\clearpage
\subsection*{ابتدای فصل دوم: از ساختار فیزیکی تا مدل ریاضی}
\begin{enumerate}
\item \textbf{بخش \secnum{۲.۱}:} ساختار \lr{Silicon}، \lr{Doping}،
\lr{Electron} و \lr{Hole}، پیوند \lr{PN} و بایاس آن؛
ساختار \lr{MOS} و حالت‌های \lr{Accumulation}، \lr{Depletion} و \lr{Inversion}؛
تعریف \lr{Threshold Voltage} و ولتاژهای پایانه‌ها.
\item \textbf{بخش \secnum{۲.۲}، مقدمات:} نقش \(W\)، \(L\)، \(t_{\mathrm{ox}}\)،
\(\mu\) و \(C_{\mathrm{ox}}\)؛ تفاوت ظرفیت و ظرفیت بر واحد سطح؛
میدان الکتریکی و سرعت حامل، با کنترل واحدها.
\item \textbf{بخش \secnum{۲.۲}، استخراج:} بار کانال، جریان \lr{Long-Channel} و
مفهوم انتگرال به‌عنوان جمع سهم‌های کوچک؛ استخراج روابط \lr{Linear} و
\lr{Saturation} و پیوستگی مرز؛ قرارداد علامت‌ها در \lr{pMOS}؛
حل مثال با تعیین ناحیه پیش و پس از محاسبه.
\item \textbf{بخش \secnum{۲.۳.۱}:} مدل ساده \lr{Gate Capacitance} و
\lr{Diffusion Capacitance}؛ اثر هندسه و اشتراک ناحیه‌های \lr{Diffusion}.
\item \textbf{بخش \secnum{۲.۳.۲}:} مدل تفصیلی ظرفیت \lr{Gate}، سهم پایانه‌ها،
\lr{Overlap} و وابستگی ظرفیت‌ها به ناحیه عملکرد.
\item \textbf{بخش \secnum{۲.۳.۳}:} مدل تفصیلی ظرفیت پیوند، سهم مساحت و محیط جانبی،
وابستگی به بایاس و مفهوم ظرفیت موثر در تغییر ولتاژ.
\end{enumerate}

\subsection*{جمع‌بندی و پیوند با دو جزوه بعدی}
در پایان جزوه، جدول نمادها و واحدها، نقشه انتخاب مدل، فرمول‌ها همراه با شرط اعتبار،
خطاهای رایج و پرسش‌های ترکیبی با پاسخ تشریحی قرار می‌گیرند.
بخش \secnum{۲.۴} به جزوه دوم و بخش \secnum{۲.۵} به جزوه سوم تعلق دارد؛ در جزوه اول فقط
ضرورت آن‌ها معرفی می‌شود. مشتق‌گیری تفصیلی برای \lr{Noise Margin} در جزوه سوم
آموزش داده خواهد شد.

\textbf{معیار پوشش:} برای هر عنوان کتاب در دامنه تعیین‌شده، توضیح آموزشی و ارجاع
قابل‌ردیابی ارائه می‌شود. منظور، بازنویسی آموزشی مفاهیم و استدلال‌هاست؛
متن کتاب یا تمام تمرین‌های پایان فصل عیناً بازتولید نمی‌شوند.
